\Needspace{18\baselineskip}
\section{Opposite specimens and the converse argument}

\begin{definition}[Opposite specimen system]
The opposite system $G^{\mathrm{op}}$ has the same specimens, modes,
entering predicate, and rooted-state predicate as $G$. It swaps the two
endpoint maps:
\[
 \ell_{G^{\mathrm{op}}}(c)=r_G(c),\qquad
 r_{G^{\mathrm{op}}}(c)=\ell_G(c).
\]
\end{definition}
\lean{LocalReasoning}{174}{187}{SpecimenSystem.opposite}
\begin{proof}[Verification of its interface laws]
Rootedness preservation is unchanged. Final-state agreement is the
reverse implication pair of the original agreement. For context-only
traversal, exchange the two rootedness hypotheses and the two alternatives
of the availability premise. Apply the original law to get a remote
specimen $d$ and two transitions. Return those transitions in the
opposite order. This explains why the interface law permits availability
on either side, rather than only on the left.
\end{proof}

\begin{lemma}[Lifting commutes with taking the opposite]
For all states $s,t$,
\[
 \mathcal L_{G^{\mathrm{op}}}(s,t)\iff\Lft(t,s).
\]
\end{lemma}
\lean{LocalReasoning}{189}{194}{SpecimenSystem.lifting_opposite}
\begin{proof}
In either direction retain the witnessing specimen and its remoteness.
Exchange the two rootedness proofs and the two endpoint equalities.
The resulting tuple is precisely the witness required by the other
lifting definition.
\end{proof}

\begin{theorem}[Converse construction of a counting simulation]
Suppose $M$ is deterministic, $(Q,P,S)$ is reasonable,
$\safe{Q^{-1}}{P}$ holds, and $G^{\mathrm{op}}$ is robust for
$(Q,P,S)$. Then $\mathcal L_{G^{\mathrm{op}}}$ is a
$(Q,P,S)$-counting simulation.
\end{theorem}
\lean{LocalReasoning}{199}{278}{SpecimenSystem.countingSimulation_converse}
The safety premise is a property of the \emph{original} system at the
converse cost relation. It is not assumed that the opposite system
already satisfies input-safety. This distinction is important when the
right execution in safety~II has length zero.

\begin{proof}
Choose a witness $c$ of $\mathcal L_{G^{\mathrm{op}}}(s,t)$, so
$s=r(c)$, $t=\ell(c)$, $c$ is remote, and both endpoints are rooted.
Finality of $r(c)$ implies finality of $\ell(c)$ by the reverse
direction of final-state agreement. This establishes condition~(A).
For condition~(B), fix a transition $r(c)\to s'$.

\emph{Settled case.} If $c$ is settled, the symmetric availability
premise of context-only traversal applies to this right transition.
Obtain a remote specimen $d$ with transitions
\[
 \ell(c)\to\ell(d),\qquad r(c)\to r(d).
\]
Determinism gives $s'=r(d)$. Propagate rootedness, and obtain
$\mathcal L_{G^{\mathrm{op}}}(r(d),\ell(d))$. Choose zero further
steps from $s'$ and one step from $\ell(c)$. The cost condition is
$1\,Q\,1$, and reflexive refinements give the up-to witness. This
covers backtracking and non-entering search.

\emph{Entering search.} Input-safety for $Q^{-1},P$ supplies three
cases, which must be read in the original endpoint order.
\begin{enumerate}[label=\textup{(\Roman*)}]
\item Use the supplied stuck execution from $r(c)$. The prescribed
  step $r(c)\to s'$ and the stuck-after-step lemma give a stuck
  execution from $s'$. The supplied stuck execution from $\ell(c)$
  is not needed in this direction.
\item There are a settled specimen $d$ and counts $k,l$ with
  \[
    \ell(c)\steps{1+k}\ell(d),\qquad r(c)\steps{l}r(d),\qquad
    l\,Q\,(1+k).
  \]
  Both endpoints of $d$ are rooted. We distinguish the two possible
  forms of $l$.

  If $l=j+1$, remove the prescribed transition from the right execution
  to get $s'\steps{j}r(d)$. Keep the left execution of length $1+k$.
  The opposite lifted pair is $(r(d),\ell(d))$, and its cost condition
  is $(1+j)\,Q\,(1+k)$, exactly the supplied relation with $l=j+1$.
  Reflexive auxiliary refinements finish this subcase.

  If $l=0$, the zero-step equality gives $r(c)=r(d)$. Thus the
  prescribed transition is an outgoing transition from $r(d)$.
  Apply context-only traversal to the \emph{settled} specimen $d$.
  Obtain a remote specimen $e$ and transitions
  \[
    \ell(d)\to\ell(e),\qquad r(d)\to r(e).
  \]
  Determinism identifies $s'$ with $r(e)$. Propagating rootedness
  gives the opposite lifted pair $(r(e),\ell(e))$. There are zero
  further source steps, whereas the target execution is
  \[
    \ell(c)\steps{1+k}\ell(d)\to\ell(e),
  \]
  of length $(1+k)+1$. Its cost obligation is
  $1\,Q\,((1+k)+1)$. The safety relation in this subcase is
  $0\,Q\,(1+k)$. Add $1\,Q\,1$ using additivity of $Q$ to
  obtain the required relation. This additional traversal is the
  essential zero-step correction; the proof never attempts to remove
  a step from a zero-step execution.
\item There are $k,l,a,b$ with
  \[
    \ell(c)\steps{1+k}a,\qquad r(c)\steps{1+l}b,\qquad
    (1+l)\,Q\,(1+k),
  \]
  and a quasi-specimen for auxiliary equivalences. Remove the prescribed
  right transition to get $s'\steps{l}b$. Both $a,b$ are rooted;
  unpacking the quasi-specimen yields a non-rewrite specimen $d$ with
  rooted endpoints and
  \[
    a\simeq_P\ell(d),\qquad r(d)\simeq_P b.
  \]
  This time take the \emph{reverse} halves:
  $b\preceq_P r(d)$ and $\ell(d)\preceq_P a$. Enlarge the
  latter auxiliary order using $P\subseteq S$ to obtain
  $\ell(d)\preceq_S a$. Since
  $\mathcal L_{G^{\mathrm{op}}}(r(d),\ell(d))$, these form
  \[
    \upto{P}{S}{\mathcal L_{G^{\mathrm{op}}}}(b,a).
  \]
  The matching counts are $l,1+k$, with the already supplied cost
  relation $(1+l)\,Q\,(1+k)$.
\end{enumerate}

\emph{Rewrite mode.} Remoteness gives $\neg\ent(c)$. Apply
robustness in $G^{\mathrm{op}}$ to the prescribed transition from
$r(c)$, which is its left endpoint. The stuck alternative is already
the required one. Otherwise the two matching executions and their cost
relation are supplied directly. Rootedness propagates along the
prescribed transition and both executions. Apply the quasi-up-to lemma
\emph{in the opposite system} to obtain the up-to relation on
$\mathcal L_{G^{\mathrm{op}}}$. This completes the proof.
\end{proof}
